\begin{tomtat}
	\begin{minipage}[t]{0.3\textwidth}
		\vspace{0pt}
		\begin{tomtat}
			\begin{center}
				\textbf{ĐỀ 32}
			\end{center}
			
		\end{tomtat}
	\end{minipage}
	\hfill
	\begin{minipage}[t]{0.5\textwidth}
		\vspace{0pt}
		\begin{tomtat}
			\begin{center}
				\includegraphics[scale=0.08]{images/32}
			\end{center}
		\end{tomtat}
	\end{minipage}
\end{tomtat}
\textbf{I. PHẦN TRẮC NGHIỆM}
\begin{ex}
	Phương trình $(x-2)(x+3)=0$ có nghiệm là?
	\choice{$x=2$}{$x=-3$}{$x=2 \text{ hoặc } x=-3$}{$x=2 \text{ và } x=-3$}
	\loigiai{
		Ta có: $(x-2)(x+3)=0$
		$\Leftrightarrow x-2=0$ hoặc $x+3=0$
		$\Leftrightarrow x=2$ hoặc $x=-3$.
		\\ \textbf{Đáp Án C}
	}
\end{ex}

\begin{ex}
	Căn bậc hai của 9 là:
	\choice{$81$}{$9$}{$3$}{$3 \text{ và } -3$}
	\loigiai{
		Căn bậc hai của một số $a \ge 0$ là một số $x$ sao cho $x^2 = a$.
		Căn bậc hai của 9 là các số $x$ sao cho $x^2 = 9$.
		Ta có $3^2 = 9$ và $(-3)^2 = 9$.
		Vậy căn bậc hai của 9 là 3 và -3.
		\\ \textbf{Đáp Án D}
	}
\end{ex}

\begin{ex}
	Nghiệm của bất phương trình $x-1 \ge 0$ là:
	\choice{$x \ge 1$}{$x > 1$}{$x=1$}{$x < 1$}
	\loigiai{
		Ta có: $x-1 \ge 0$
		$\Leftrightarrow x \ge 1$.
		\\ \textbf{Đáp Án A}
	}
\end{ex}

\begin{ex}
	Điểm nào sau đây thuộc đồ thị hàm số $y=2x^2$?
	\choice{$(0; 2)$}{$(1; 2)$}{$(-1; -2)$}{$(1; 0)$}
	\loigiai{
		Xét các điểm:
		A. Với $(0; 2)$: $y = 2(0)^2 = 0 \neq 2$. Loại.
		B. Với $(1; 2)$: $y = 2(1)^2 = 2 \cdot 1 = 2$. Chọn.
		C. Với $(-1; -2)$: $y = 2(-1)^2 = 2 \cdot 1 = 2 \neq -2$. Loại.
		D. Với $(1; 0)$: $y = 2(1)^2 = 2 \neq 0$. Loại.
		Vậy điểm $(1;2)$ thuộc đồ thị hàm số $y=2x^2$.
		\\ \textbf{Đáp Án B}
	}
\end{ex}

\begin{ex}
	Cho tam giác ABC vuông tại A có $\widehat{B}=\alpha$.
	\choice{$\sin \alpha = \frac{AB}{AC}$}{$\sin \alpha = \frac{AC}{AB}$}{$\sin \alpha = \frac{AC}{BC}$}{$\sin \alpha = \frac{AB}{BC}$}
	\loigiai{
		Trong tam giác ABC vuông tại A, với góc $\widehat{B}=\alpha$:
		Cạnh đối của góc $\alpha$ là AC.
		Cạnh kề của góc $\alpha$ là AB.
		Cạnh huyền là BC.
		Theo định nghĩa tỉ số lượng giác sin:
		$\sin \alpha = \frac{\text{cạnh đối}}{\text{cạnh huyền}} = \frac{AC}{BC}$.
		\\ \textbf{Đáp Án C}
	}
\end{ex}
\begin{ex}
	Cho hình nón có bán kính đáy là 2cm, đường sinh là 4cm. Hỏi diện tích xung quanh là bao nhiêu?
	\choice{$5\pi \text{ cm}^2$}{$6\pi \text{ cm}^2$}{$7\pi \text{ cm}^2$}{$8\pi \text{ cm}^2$}
	\loigiai{
		Bán kính đáy của hình nón là $r = 2$ cm.
		Độ dài đường sinh của hình nón là $l = 4$ cm.
		Diện tích xung quanh của hình nón được tính bằng công thức $S_{xq} = \pi r l$.
		$S_{xq} = \pi \cdot 2 \cdot 4 = 8\pi \text{ (cm}^2\text{)}$.
		\\ \textbf{Đáp Án D}
	}
\end{ex}

\begin{ex}
	Gieo một con xúc sắc 50 lần và được kết quả như sau:
	\begin{center}
		\begin{tabular}{|c|c|c|c|c|c|c|}
			\hline
			Số chấm xuất hiện & 1 & 2 & 3 & 4 & 5 & 6 \\
			\hline
			Tần số & 8 & 7 & 9 & 8 & 6 & 12 \\
			\hline
		\end{tabular}
	\end{center}
	Tần số xuất hiện mặt ba chấm là:
	\choice{$8$}{$9$}{$10$}{$11$}
	\loigiai{
		Dựa vào bảng tần số đã cho, tần số xuất hiện mặt ba chấm là 9.
		\\ \textbf{Đáp Án B}
	}
\end{ex}

\begin{ex}
	Một hộp chứa 4 quả cầu cùng loại trong đó có 1 quả cầu đỏ, 1 quả cầu xanh và 2 quả cầu vàng. Chọn ngẫu nhiên 1 quả cầu. Xác suất của biến cố “Chọn được 1 quả màu vàng” là:
	\choice{$\frac{1}{6}$}{$\frac{1}{4}$}{$\frac{1}{2}$}{$\frac{2}{3}$}
	\loigiai{
		Tổng số quả cầu trong hộp là $1 (\text{đỏ}) + 1 (\text{xanh}) + 2 (\text{vàng}) = 4$ quả cầu.
		Số phần tử của không gian mẫu là $n(\Omega) = 4$.
		Gọi A là biến cố “Chọn được 1 quả màu vàng”.
		Số quả cầu màu vàng là 2.
		Số kết quả thuận lợi cho biến cố A là $n(A) = 2$.
		Xác suất của biến cố A là: $P(A) = \frac{n(A)}{n(\Omega)} = \frac{2}{4} = \frac{1}{2}$.
		\\ \textbf{Đáp Án C}
	}
\end{ex}

\newpage
\textbf{II. PHẦN TỰ LUẬN}
\begin{ex}
	Rút gọn biểu thức A biết $A = \frac{\sqrt{x}+2}{\sqrt{x}-2} \cdot \left( \frac{x+4}{x-4} - \frac{2}{\sqrt{x}-2} \right)$, với $x \ge 0, x \neq 4$.
	\loigiai{
		Điều kiện: $x \ge 0, x \neq 4$.
		Ta có $x-4 = (\sqrt{x}-2)(\sqrt{x}+2)$.
		\begin{align*} A &= \frac{\sqrt{x}+2}{\sqrt{x}-2} \cdot \left( \frac{x+4}{(\sqrt{x}-2)(\sqrt{x}+2)} - \frac{2(\sqrt{x}+2)}{(\sqrt{x}-2)(\sqrt{x}+2)} \right) \\ &= \frac{\sqrt{x}+2}{\sqrt{x}-2} \cdot \frac{x+4 - 2\sqrt{x}-4}{(\sqrt{x}-2)(\sqrt{x}+2)} \\ &= \frac{\sqrt{x}+2}{\sqrt{x}-2} \cdot \frac{x - 2\sqrt{x}}{(\sqrt{x}-2)(\sqrt{x}+2)} \\ &= \frac{\sqrt{x}+2}{\sqrt{x}-2} \cdot \frac{\sqrt{x}(\sqrt{x}-2)}{(\sqrt{x}-2)(\sqrt{x}+2)} \end{align*}
		Vì $x \ge 0, x \neq 4$, nên $\sqrt{x}-2 \neq 0$ và $\sqrt{x}+2 \neq 0$.
		\begin{align*} A &= \frac{\cancel{(\sqrt{x}+2)}}{\cancel{(\sqrt{x}-2)}} \cdot \frac{\sqrt{x}\cancel{(\sqrt{x}-2)}}{(\sqrt{x}-2)\cancel{(\sqrt{x}+2)}} \\ &= \frac{\sqrt{x}}{\sqrt{x}-2} \end{align*}
		Vậy $A = \frac{\sqrt{x}}{\sqrt{x}-2}$.
	}
\end{ex}

\begin{ex}
	Giải hệ phương trình $\begin{cases} 2x+y=3 \\ x-y=6 \end{cases}$.
	\loigiai{
		Ta có hệ phương trình:
		$$ \begin{cases} 2x+y=3 \quad &(1) \\ x-y=6 \quad &(2) \end{cases} $$
		Cộng vế theo vế của phương trình (1) và (2), ta được:
		\begin{align*} (2x+y) + (x-y) &= 3+6 \\ 3x &= 9 \\ x &= \frac{9}{3} \\ x &= 3 \end{align*}
		Thay $x=3$ vào phương trình (2):
		\begin{align*} 3-y &= 6 \\ y &= 3-6 \\ y &= -3 \end{align*}
		Vậy nghiệm của hệ phương trình là $(x;y) = (3;-3)$.
	}
\end{ex}

\begin{ex}
	Biết phương trình $x^2+x-6=0$ luôn có hai nghiệm phân biệt $x_1; x_2$. Không cần giải phương trình hãy tính giá trị của biểu thức $A = x_1^2+x_2^2+x_1x_2$.
	\loigiai{
		Phương trình $x^2+x-6=0$ có $a=1, b=1, c=-6$.
		$\Delta = b^2-4ac = 1^2 - 4(1)(-6) = 1+24=25 > 0$.
		Phương trình có hai nghiệm phân biệt $x_1, x_2$.
		Theo hệ thức Vi-ét, ta có:
		$$ \begin{cases} x_1 + x_2 = -\frac{b}{a} = -\frac{1}{1} = -1 \\ x_1 x_2 = \frac{c}{a} = \frac{-6}{1} = -6 \end{cases} $$
		Ta biến đổi biểu thức A:
		\begin{align*} A &= x_1^2+x_2^2+x_1x_2 \\ &= (x_1+x_2)^2 - 2x_1x_2 + x_1x_2 \\ &= (x_1+x_2)^2 - x_1x_2 \end{align*}
		Thay giá trị từ hệ thức Vi-ét vào biểu thức A:
		\begin{align*} A &= (-1)^2 - (-6) \\ &= 1 + 6 \\ &= 7 \end{align*}
		Vậy $A=7$.
	}
\end{ex}
\begin{ex}
	Viết ngẫu nhiên một số tự nhiên có hai chữ số nhỏ hơn 100. Có bao nhiêu kết quả có thể xảy ra của phép thử trên? Tính xác suất của mỗi biến cố sau:
	B: “Số tự nhiên được viết ra là số chẵn”;
	\loigiai{
		Các số tự nhiên có hai chữ số nhỏ hơn 100 là các số từ 10 đến 99.
		Số các kết quả có thể xảy ra của phép thử (số phần tử của không gian mẫu $\Omega$) là:
		$n(\Omega) = 99 - 10 + 1 = 90$.
		
		Biến cố B: “Số tự nhiên được viết ra là số chẵn”.
		Các số tự nhiên chẵn có hai chữ số nhỏ hơn 100 là: $10; 12; 14; \dots ; 96; 98$.
		Để tìm số các số chẵn này, ta có thể dùng công thức tính số số hạng của một dãy số cách đều:
		Số số hạng = (Số cuối - Số đầu) / Khoảng cách + 1.
		Số các số chẵn là: $n(B) = \frac{98 - 10}{2} + 1 = \frac{88}{2} + 1 = 44 + 1 = 45$.
		
		Xác suất của biến cố B là:
		$P(B) = \frac{n(B)}{n(\Omega)} = \frac{45}{90} = \frac{1}{2}$.
		
		\textbf{Kết luận:}
		\begin{itemize}
			\item Số kết quả có thể xảy ra của phép thử là 90.
			\item Xác suất của biến cố B: “Số tự nhiên được viết ra là số chẵn” là $\frac{1}{2}$.
		\end{itemize}
	}
\end{ex}
\begin{ex}
	Một xí nghiệp dự định làm sân bóng nhân tạo 5 người có hình chữ nhật với chu vi bằng 110m. Người ta tính rằng có thể chuyển thành sân bóng nhân tạo 7 người bằng cách tăng chiều rộng lên 2,5 lần và chiều dài lên gấp đôi, khi đó chu vi của sân bóng mới là 240m. Hãy tính diện tích của sân bóng dự định làm ban đầu.
	\loigiai{
		Gọi chiều dài ban đầu của sân bóng là $x$ (mét).
		Gọi chiều rộng ban đầu của sân bóng là $y$ (mét).
		Điều kiện: $x > 0, y > 0$.
		
		Chu vi ban đầu của sân bóng là 110m, nên ta có phương trình:
		\begin{align*} 2(x+y) &= 110 \\ x+y &= 55 \quad (1) \end{align*}
		
		Khi chuyển thành sân bóng 7 người:
		Chiều dài mới là $2x$ (mét).
		Chiều rộng mới là $2,5y$ (mét).
		Chu vi của sân bóng mới là 240m, nên ta có phương trình:
		\begin{align*} 2(2x + 2,5y) &= 240 \\ 2x + 2,5y &= 120 \quad (2) \end{align*}
		
		Từ (1) và (2), ta có hệ phương trình:
		$$ \begin{cases} x + y = 55 \\ 2x + 2,5y = 120 \end{cases} $$
		Giải hệ phương trình:
		Từ (1) $\Rightarrow x = 55 - y$.
		Thế $x = 55 - y$ vào (2):
		\begin{align*} 2(55 - y) + 2,5y &= 120 \\ 110 - 2y + 2,5y &= 120 \\ 0,5y &= 120 - 110 \\ 0,5y &= 10 \\ y &= \frac{10}{0,5} \\ y &= 20 \end{align*}
		Thay $y = 20$ vào $x = 55 - y$:
		\begin{align*} x &= 55 - 20 \\ x &= 35 \end{align*}
		(Thỏa mãn điều kiện $x > 0, y > 0$).
		
		Vậy, chiều dài ban đầu của sân bóng là 35 mét và chiều rộng ban đầu là 20 mét.
		Diện tích của sân bóng dự định làm ban đầu là:
		$S = x \cdot y = 35 \cdot 20 = 700 \text{ (m}^2\text{)}$.
	}
\end{ex}
\begin{ex}
	Cho tam giác ABC nhọn. Đường tròn (O) đường kính BC cắt AB, AC lần lượt tại E và D; BD cắt CE tại H, AH cắt BC tại I. Từ A kẻ tiếp tuyến AM, AN của đường tròn (O) (M, N là tiếp điểm).
	\begin{enumerate}
		\item[a)] Chứng minh tứ giác AEHD nội tiếp.
		\item[b)] Chứng minh $AB \cdot BE = BI \cdot BC$ từ đó suy ra $AB \cdot BE + AC \cdot CD = BC^2$.
	\end{enumerate}
\end{ex}

\begin{ex}
	\begin{enumerate}
		\item[a)] Một bồn nước I-nox có dạng hình trụ với chiều cao 2m và chu vi đáy là $\pi$ (m). Hỏi bồn nước này đựng đầy được bao nhiêu mét khối nước? (Bỏ qua bề dày của vỏ bồn nước và lấy $\pi \approx 3,14$).
		\item[b)] Cho các số $a, b, c > 0$ và $a+b+c \le \frac{3}{2}$. Tìm giá trị nhỏ nhất của biểu thức:
		$$ Q = \sqrt{a^2+\frac{1}{b^2}} + \sqrt{b^2+\frac{1}{c^2}} + \sqrt{c^2+\frac{1}{a^2}} $$
	\end{enumerate}
\end{ex}

\setcounterref{ex}{0}
